\documentclass[journal,onecolumn]{IEEEtran}

\usepackage{amsmath,amssymb}
\usepackage{array}
\usepackage{booktabs}
\usepackage{cite}
\usepackage{float}
\usepackage{makecell}
\usepackage{multirow}
\usepackage{newtxtext,newtxmath}
\usepackage{siunitx}
\usepackage[hidelinks]{hyperref}

\hypersetup{
  pdftitle={Supplementary Material for PA-LOI: A Proactive Risk Layer for Occluded Pedestrian Emergence},
  pdfauthor={Hongyu Peng; Shangqi Zhang; Bo Chen; Xiao Wang}
}

\sisetup{detect-all=true}
\renewcommand{\thetable}{S\arabic{table}}

\title{Supplementary Material for ``PA-LOI: A Proactive Risk Layer for Occluded Pedestrian Emergence''}
\author{Hongyu~Peng, Shangqi~Zhang, Bo~Chen, and Xiao~Wang}

\begin{document}
\maketitle

\section{Literature Positioning}

Table~\ref{tab:related-comparison} complements the main-paper discussion by
comparing hidden-risk representation, planner coupling, evaluation scale, and
the availability of fixed hazard geometry. The entries are qualitative because
the cited systems use different planners and scenario definitions.

\begin{table*}[!ht]
\caption{Qualitative Positioning of PA-LOI Among Representative Occlusion-Aware Approaches. The Comparison Concerns Design Choices, Not Head-to-Head Metrics, Because the Cited Methods Are Coupled to Different Planner Stacks and Scenario Definitions.}
\label{tab:related-comparison}
\centering
\scriptsize
\setlength{\tabcolsep}{4pt}
\begin{tabular}{>{\raggedright\arraybackslash}p{0.20\textwidth}>{\raggedright\arraybackslash}p{0.23\textwidth}>{\raggedright\arraybackslash}p{0.23\textwidth}>{\raggedright\arraybackslash}p{0.13\textwidth}>{\raggedright\arraybackslash}p{0.10\textwidth}}
\toprule
Approach & Hidden-risk representation & Planner coupling & Closed-loop evaluation scale & Fixed-geometry protocol \\
\midrule
Set-based reachability~\cite{orzechowski2018tackling,koschi2021setbased} & Reachable occupancy behind occluders and beyond sensor range & Verification or constraint layer over a given planner & Formal analysis with case studies & Not released \\
Dynamic occlusion shadows~\cite{nager2019shadows} & Tracked shadow regions with potential occupants & Reactive constraints on speed and clearance & Closed-loop case studies & Not released \\
POMDP belief planning~\cite{zhang2021pomdp} & Belief over hidden-agent states & Belief-space planner replaces the nominal stack & Simulation case studies & Not released \\
Occlusion-aware contingency planning~\cite{zheng2026occlusioncontingency} & Hypothetical agents with contingency branches & Dedicated safety-critical MPC formulation & Simulation scenario studies & Not released \\
Visibility maximization~\cite{narksri2022visibility} & Visibility deficit along candidate paths & Modified objective of the nominal planner & Intersection case studies & Not released \\
\midrule
PA-LOI (this work) & Conflict-centered phantom source with safe-speed bound from bounded-rational emergence & Additive velocity-hinge cost in the trajectory optimizer; learned prediction and scenario-tree logic preserved; AEB fallback & 4050 closed-loop AV2 replay runs: five systems on 81 scenes at six trigger distances, plus no-ghost, walking-speed, extended-distance, and parameter protocols & JSONL-fixed geometry; archive supplied (records, code, scripts) \\
\bottomrule
\end{tabular}
\end{table*}


\section{Extended Trigger-Distance Results}

Table~\ref{tab:extended-distance} reports every distance--speed cell from the
extended moving-pedestrian protocol summarized in Section~V-E of the main
paper. The original per-run records are included in the accompanying
reproducibility archive.

\begin{table*}[!ht]
\caption{Moving-Pedestrian Emergence at Extended Trigger Distances (27 Scenes, 800 Frames). $v_{\mathrm{imp}}^2$ Columns Are Hit-Only Means, as in the Main-Paper Pedestrian-Speed Table. Best Low-Severity and Impact-Energy Values per Distance--Speed Cell in Bold.}
\label{tab:extended-distance}
\centering
\scriptsize
\setlength{\tabcolsep}{4.5pt}
\begin{tabular}{clcccccc}
\toprule
 & & \multicolumn{3}{c}{$d_{\mathrm{trig}}=\SI{4.5}{m}$} & \multicolumn{3}{c}{$d_{\mathrm{trig}}=\SI{3.5}{m}$} \\
\cmidrule(lr){3-5}\cmidrule(lr){6-8}
\makecell{Speed\\(m/s)} & Method & Low-severity & Zero coll. & \makecell{Mean $v_{\mathrm{imp}}^2$\\(m$^2$/s$^2$)} & Low-severity & Zero coll. & \makecell{Mean $v_{\mathrm{imp}}^2$\\(m$^2$/s$^2$)} \\
\midrule
\multirow{3}{*}{1.0} & PA-LOI + AEB & \textbf{26/27 (96.3\%)} & 2/27 (7.4\%) & \textbf{1.00} & \textbf{26/27 (96.3\%)} & 0/27 (0.0\%) & \textbf{0.67} \\
 & AEB-only & 21/27 (77.8\%) & 1/27 (3.7\%) & 4.36 & 20/27 (74.1\%) & 0/27 (0.0\%) & 4.83 \\
 & MIND & 5/27 (18.5\%) & 4/27 (14.8\%) & 25.41 & 6/27 (22.2\%) & 5/27 (18.5\%) & 25.01 \\
\midrule
\multirow{3}{*}{2.0} & PA-LOI + AEB & \textbf{26/27 (96.3\%)} & 5/27 (18.5\%) & \textbf{2.78} & \textbf{26/27 (96.3\%)} & 0/27 (0.0\%) & \textbf{3.01} \\
 & AEB-only & 14/27 (51.9\%) & 1/27 (3.7\%) & 11.95 & 13/27 (48.1\%) & 1/27 (3.7\%) & 12.81 \\
 & MIND & 3/27 (11.1\%) & 0/27 (0.0\%) & 28.05 & 3/27 (11.1\%) & 0/27 (0.0\%) & 27.25 \\
\midrule
\multirow{3}{*}{3.0} & PA-LOI + AEB & \textbf{26/27 (96.3\%)} & 4/27 (14.8\%) & \textbf{3.15} & \textbf{25/27 (92.6\%)} & 0/27 (0.0\%) & \textbf{4.80} \\
 & AEB-only & 12/27 (44.4\%) & 1/27 (3.7\%) & 15.18 & 10/27 (37.0\%) & 0/27 (0.0\%) & 16.41 \\
 & MIND & 3/27 (11.1\%) & 0/27 (0.0\%) & 27.81 & 4/27 (14.8\%) & 0/27 (0.0\%) & 27.40 \\
\bottomrule
\end{tabular}

\vspace{2pt}{\scriptsize In one scene AEB-only hits replay background traffic before the trigger in all six cells (counted unsafe); excluding that scene for all stacks changes no rate by more than 3.1\,pp and no ordering.}
\end{table*}


\section{Risk-Layer Parameter Perturbations}

Table~\ref{tab:param-ablation} reports the six one-dimensional perturbations
summarized in Section~V-F of the main paper. The deployed setting is repeated
to make the safety--efficiency changes directly auditable.

\begin{table*}[!ht]
\caption{Risk-Layer Parameter Perturbation on the 27-Scene Set (Ghost at 5.5 and 4.0\,m, Instant-Center, 650 Frames, Plus No-Ghost Efficiency). The Deployed Row Uses the Published Strict-650 Tables on the Identical Cells.}
\label{tab:param-ablation}
\centering
\scriptsize
\setlength{\tabcolsep}{6pt}
\begin{tabular}{lccccccc}
\toprule
Setting & Spawn & \makecell{Low-severity\\(all 54 cells)} & \makecell{Low-severity\\(spawned)} & \makecell{Zero collision\\(all 54 cells)} & \makecell{Mean $v_{\mathrm{imp}}^2$\\(spawned)} & \makecell{No-ghost\\speed (m/s)} & \makecell{PF/RC} \\
\midrule
Deployed ($0.6/2.0/25$) & 54/54 & 92.6\% & 92.6\% & 48.1\% & 2.87 & 3.47 & 0/0 \\
\midrule
$\eta=0.4$ & 50/54 & 87.0\% & 90.0\% & 44.4\% & 2.55 & 3.29 & 0/3 \\
$\eta=0.8$ & 54/54 & 87.0\% & 87.0\% & 38.9\% & 3.24 & 3.57 & 0/4 \\
$v_{\mathrm{floor}}=1.0$ & 51/54 & 88.9\% & 90.2\% & 46.3\% & 2.72 & 3.34 & 0/3 \\
$v_{\mathrm{floor}}=3.0$ & 51/54 & 94.4\% & 94.1\% & 50.0\% & 2.86 & 3.55 & 0/1 \\
$w_{\max}=15$ & 54/54 & 81.5\% & 81.5\% & 37.0\% & 3.34 & 3.54 & 1/4 \\
$w_{\max}=35$ & 51/54 & 88.9\% & 90.2\% & 44.4\% & 2.44 & 3.31 & 0/3 \\
\bottomrule
\end{tabular}

\vspace{2pt}{\scriptsize Unspawned cells (ego never reaches the trigger within 650 frames) count as safe in all-cell columns; the spawned-only column removes this credit. PF/RC: planner failures / replay collisions over 27 no-ghost runs; the $w_{\max}=15$ speed includes one run truncated at frame 630 (3.46\,m/s excluding it).}
\end{table*}


\section{Unfiltered Candidate-Pool Audit}

The AV2 search returned 90 new geometric candidates. Two lacked complete
three-stack screening records and were not used; 60 of the remaining 88 form
new60, while 28 were screened out by the protocol-integrity checks described
in Section~IV-B. Table~\ref{tab:candidate-pool-audit} removes this clean-scene
filter and retains all 88 log-complete candidates. The strict 650-frame
evaluable-set rule is applied exactly as in the main paper, background
collisions remain unsafe, and non-spawning is exposed by the spawn rate rather
than credited as safety. This is intentionally a harsher, confounded pool: the
absolute rates fall, but PA-LOI's ordering over AEB-only and MIND remains. The
scene-clustered 95\% intervals for low severity are $[43.0,61.4]$, $[9.1,22.8]$,
and $[1.0,9.2]$, respectively. The archive gives all 90 membership decisions
and task-level screening flags; no mutually exclusive exclusion reasons are
assigned retrospectively because flags can overlap.

% AUTO-GENERATED by compute_candidate_pool_audit.py.
\begin{table}[H]
\caption{Unfiltered New90 Candidate-Pool Audit. All 88 Geometric Candidates With Complete Strict-650 Records Are Retained, Including the 28 Scenes Excluded From New60. Safety Is Computed on Evaluable Tasks; Spawn Rate Exposes Non-Encounter Conservatism.}
\label{tab:candidate-pool-audit}
\centering
\scriptsize
\setlength{\tabcolsep}{7pt}
\begin{tabular}{lccccccc}
\toprule
Method & Tasks & Spawn rate & Eval. tasks & Low-severity & Zero collision & Background collisions & Mean $v_{\mathrm{imp}}^2$ \\
\midrule
PA-LOI + AEB & 528 & 73.9\% & 348 & 182/348 (52.3\%) & 47/348 (13.5\%) & 63 & 10.15 \\
AEB-only & 528 & 92.4\% & 471 & 74/471 (15.7\%) & 12/471 (2.5\%) & 15 & 31.96 \\
MIND & 528 & 98.1\% & 510 & 23/510 (4.5\%) & 0/510 (0.0\%) & 31 & 42.06 \\
\bottomrule
\end{tabular}
\end{table}


\section{Conference-to-Journal Extension}

Table~\ref{tab:conference-delta} makes the extension beyond the accepted ITSC
paper explicit~\cite{paloi_itsc2026}. The risk-layer algorithm is not presented
as new in this journal submission; the new contribution is the question and
evidence enabled by clean81, the five-system comparison, exposure accounting,
robustness protocols, statistical analysis, and auditable release.

\begin{table}[H]
\caption{Itemized Difference Between the ITSC Conference Paper and This Journal Article. The Risk-Layer Algorithm Is Explicitly Carried Over; the Journal Contribution Is the New Scientific Question, Evaluation Instrument, Comparative Evidence, and Reproducibility Record.}
\label{tab:conference-delta}
\centering
\scriptsize
\setlength{\tabcolsep}{4pt}
\renewcommand{\arraystretch}{1.12}
\begin{tabular}{>{\raggedright\arraybackslash}p{0.15\textwidth}>{\raggedright\arraybackslash}p{0.36\textwidth}>{\raggedright\arraybackslash}p{0.41\textwidth}}
\toprule
Aspect & ITSC conference paper & Present journal article \\
\midrule
Scientific question & Demonstrates whether proactive occlusion-risk shaping can improve on late AEB in representative scenes. & Tests whether the unchanged layer reduces severity with less conservatism than strict occlusion bounds on identical hazards, and quantifies the no-pedestrian cost. \\
Risk layer & Introduces TTA gating, the crossing-time safe-speed construction, velocity-hinge cost, and AEB fallback. & Preserves those components; adds an implementation-faithful formalization, complete deployed parameters and fallback semantics, and the bounded-rational policy-class interpretation. \\
Evaluation scale & Two representative AV2-derived occlusion sites with trigger-distance sweeps and an illustrative threat-free trial. & clean81 with 81 fixed scenes, six trigger distances, five stacks, and 4050 closed-loop runs across the main, no-ghost, walking, extended-distance, and parameter protocols. \\
Comparators & MIND, AEB-only, PA-LOI, and a generic conservative worst-case assumption. & MIND, AEB-only, PA-LOI, and traceable per-step ports of set-based reachability and dynamic-shadow families under shared geometry, physics, and AEB shielding. \\
Reporting and inference & Collision/impact summaries, one threat-free example, component examples, and module latency. & Low-severity and zero collision separated; background collision, planner infeasibility, impact speed, velocity-squared proxy, spawn-rate conservatism, scene-clustered intervals, paired effect sizes, and exact sign tests. \\
Robustness and audit & Representative component ablations and qualitative limitations. & Held-out new60 split, thresholds from 2.0--3.5\,m/s, moving pedestrians over three triggers, six parameter perturbations, both horizon accountings, and an unfiltered 88-candidate curation audit. \\
Reproducibility & Scenario-level implementation demonstration. & Canonical JSONL geometry, archived per-run records, method and supervisor code, deterministic table scripts, checksums, and a tested workspace-restoration map. \\
\bottomrule
\end{tabular}
\end{table}


\bibliographystyle{IEEEtran}
\bibliography{references}

\end{document}
